% ════════════════════════════════════════════════════════════════
%  CHAPITRE 5 — TESTS ET VALIDATION
% ════════════════════════════════════════════════════════════════

\chapter{Tests et validation}
\graphicspath{{Isipfe/Figures/}}

% ── Introduction du chapitre ────────────────────────────────────
\section*{Introduction}
\addcontentsline{toc}{section}{Introduction}

Ce chapitre présente la démarche de test appliquée à la plateforme :
stratégie, environnement, cas de test, résultats et validation finale du
système.

L'objectif est de vérifier que la solution est correcte, stable et conforme
aux besoins exprimés dans les chapitres précédents.


% ════════════════════════════════════════════════════════════════
\section{Stratégie de test}
% ════════════════════════════════════════════════════════════════

La stratégie de test suit une logique progressive : tester d'abord les
composants de base, puis les interactions entre modules, et enfin le
comportement global côté utilisateur.

\textbf{Types de tests retenus :}
\begin{itemize}
    \item \textbf{Tests unitaires} : validation automatisée de fonctions
    ciblées (principalement côté interface RH, via Vitest).
    \item \textbf{Tests d'intégration} : vérification des échanges entre
    API, base de données, notifications et services externes.
    \item \textbf{Tests fonctionnels} : validation des scénarios métier
    complets (candidat, RH, administrateur).
\end{itemize}

Cette combinaison réduit le risque de régression et garantit que le
comportement final reste cohérent avec les exigences métier.

% ════════════════════════════════════════════════════════════════
\section{Environnement de test}
% ════════════════════════════════════════════════════════════════

\subsection{Outils}

Les outils mobilisés pour la phase de test sont les suivants :
\begin{itemize}
    \item \textbf{curl} : tests des endpoints REST et capture des codes
    HTTP pour les preuves.
    \item \textbf{Vitest} : tests unitaires cibles sur l'interface RH.
    \item \textbf{Navigateur (Chrome/Firefox)} : vérification des parcours
    fonctionnels sur interfaces web.
    \item \textbf{Android (PWA)} : validation du comportement mobile et de
    l'installation de l'application.
    \item \textbf{Prisma Studio / PostgreSQL} : contrôle des données
    persistées après exécution des scénarios.
\end{itemize}

\textit{Remarque :} l'automatisation est volontairement ciblée (Vitest),
tandis que les validations fonctionnelles et d'intégration sont effectuées
pas à pas via curl et les interfaces.

\subsection{Configuration}

Les tests ont été exécutés dans un environnement local proche des conditions
réelles :
\begin{itemize}
    \item backend Node.js/Express connecté à PostgreSQL ;
    \item interface candidat (PWA), interface RH et interface Admin ;
    \item communication temps réel via Socket.io ;
    \item envoi email activé pour les scénarios d'entretien.
\end{itemize}

Cette configuration a permis de vérifier les flux de bout en bout, sans
simuler artificiellement les interactions principales du système.

% ════════════════════════════════════════════════════════════════
\section{Conception des cas de test}
% ════════════════════════════════════════════════════════════════

\subsection{Méthodes (équivalence, limites)}

La conception des cas de test repose sur deux méthodes simples et efficaces :
\begin{itemize}
    \item \textbf{Partition d'équivalence} : regroupement des entrées
    valides et invalides (ex. formats email, types de fichiers CV).
    \item \textbf{Analyse des valeurs limites} : vérification des bornes
    (ex. nombre maximal de tentatives, transitions d'état terminales).
\end{itemize}

Ces méthodes permettent de couvrir les scénarios critiques sans multiplier
inutilement les cas de test.

\subsection{Tableaux de test}

\begin{table}[H]
\centering
\caption{Cas de test --- authentification et candidature}
\begin{tabularx}{\textwidth}{|X|X|l|}
\hline
\textbf{Cas de test} & \textbf{Résultat attendu} & \textbf{Statut} \\
\hline
Inscription avec email invalide & Rejet avec message explicite &
\textcolor{green!60!black}{OK} \\
\hline
Connexion avant vérification email & Accès refusé &
\textcolor{green!60!black}{OK} \\
\hline
Upload CV non PDF & Fichier rejeté avant stockage &
\textcolor{green!60!black}{OK} \\
\hline
Candidature en doublon (même offre) & Réponse \texttt{409} &
\textcolor{green!60!black}{OK} \\
\hline
\end{tabularx}
\end{table}

\begin{table}[H]
\centering
\caption{Cas de test --- RH, workflow et sécurité}
\begin{tabularx}{\textwidth}{|X|X|l|}
\hline
\textbf{Cas de test} & \textbf{Résultat attendu} & \textbf{Statut} \\
\hline
Changement de statut depuis Kanban & Mise à jour en base + notification &
\textcolor{green!60!black}{OK} \\
\hline
Planification d'entretien sur dossier finalisé & Opération bloquée
(\texttt{409}) & \textcolor{green!60!black}{OK} \\
\hline
Accès RH à un dossier d'un autre site & Accès interdit (\texttt{403}) &
\textcolor{green!60!black}{OK} \\
\hline
Route admin appelée par un RH & Accès interdit (\texttt{403}) &
\textcolor{green!60!black}{OK} \\
\hline
\end{tabularx}
\end{table}

% ════════════════════════════════════════════════════════════════
\section{Exécution des tests}
% ════════════════════════════════════════════════════════════════

\subsection{Résultats}

L'exécution des tests montre que les scénarios métiers principaux sont
validés. Les flux candidat, RH et administrateur se comportent
conformément aux attentes dans les cas testés.

Les vérifications de sécurité (rôles, isolation par site, accès ressources)
ont également donné des résultats satisfaisants.

Les tests unitaires automatisés sont présents de manière ciblée sur des
modules RH (authentification locale, analyse duale). Les appels curl
confirment les retours attendus (200/401/403/409) sur les cas critiques.

\subsection{Analyse}

L'analyse des résultats confirme trois points importants :
\begin{itemize}
    \item la plateforme reste cohérente même en présence d'entrées
    invalides ;
    \item les règles métier critiques sont correctement appliquées ;
    \item la traçabilité des actions facilite la détection d'écarts.
\end{itemize}

Les scénarios négatifs validés (accès interdit, doublon, vérification email)
montrent que le système applique des contraintes d'ingénierie, et pas
seulement des parcours nominaux.

Quelques ajustements mineurs d'interface ont été observés pendant les tests
(formatage d'affichage, messages utilisateur), puis corrigés dans la même
itération.

\subsection{Traces d'exécution}

Les captures suivantes proviennent d'appels curl (codes HTTP) et de Vitest.
Elles couvrent la santé API, le contrôle d'accès, les doublons et la
gouvernance admin, avec des résultats observables.

\begin{figure}[H]
\centering
\includegraphics[width=0.90\textwidth]{proof1.jpg}
\caption{Santé de l'API (HTTP 200)}
\end{figure}

\begin{figure}[H]
\centering
\includegraphics[width=0.90\textwidth]{proof2.jpg}
\caption{Route protégée sans jeton (HTTP 401)}
\end{figure}

\begin{figure}[H]
\centering
\includegraphics[width=0.90\textwidth]{proof3.jpg}
\caption{Authentification candidat (HTTP 200)}
\end{figure}

\begin{figure}[H]
\centering
\includegraphics[width=0.90\textwidth]{proof4.jpg}
\caption{Accès admin interdit au candidat (HTTP 403)}
\end{figure}

\begin{figure}[H]
\centering
\includegraphics[width=0.90\textwidth]{proof5.jpg}
\caption{Candidature en doublon (HTTP 409)}
\end{figure}

\begin{figure}[H]
\centering
\includegraphics[width=0.90\textwidth]{proof6.jpg}
\caption{E-mail déjà enregistré (HTTP 409)}
\end{figure}

\begin{figure}[H]
\centering
\includegraphics[width=0.90\textwidth]{proof7.jpg}
\caption{Accès RH par site (HTTP 200)}
\end{figure}

\begin{figure}[H]
\centering
\includegraphics[width=0.90\textwidth]{proof8.jpg}
\caption{Gouvernance admin (HTTP 200)}
\end{figure}

\begin{figure}[H]
\centering
\includegraphics[width=0.90\textwidth]{vitest.jpg}
\caption{Tests unitaires (Vitest)}
\end{figure}

% ════════════════════════════════════════════════════════════════
\section{Validation du système}
% ════════════════════════════════════════════════════════════════

\subsection{Vérification des exigences}

Le tableau ci-dessous synthétise la correspondance entre exigences
fonctionnelles et fonctionnalités livrées.

\begin{table}[H]
\centering
\caption{Vérification des exigences fonctionnelles}
\begin{tabularx}{\textwidth}{|X|l|}
\hline
\textbf{Exigence} & \textbf{Validation} \\
\hline
Consulter les offres depuis mobile & \textcolor{green!60!black}{Validée} \\
\hline
Soumettre une candidature avec CV & \textcolor{green!60!black}{Validée} \\
\hline
Suivre l'évolution de candidature & \textcolor{green!60!black}{Validée} \\
\hline
Traiter les dossiers via pipeline RH & \textcolor{green!60!black}{Validée} \\
\hline
Planifier un entretien avec notification & \textcolor{green!60!black}{Validée} \\
\hline
Gérer les comptes RH côté admin & \textcolor{green!60!black}{Validée} \\
\hline
\end{tabularx}
\end{table}

\subsection{Traçabilité exigences-tests}

Pour renforcer la validation, le tableau suivant relie chaque exigence
à la méthode de test appliquée et à l'élément vérifié.

\begin{table}[H]
\centering
\caption{Traçabilité entre exigences et vérifications}
\begin{tabularx}{\textwidth}{|X|l|X|}
\hline
\textbf{Exigence} & \textbf{Type de test} & \textbf{Élément de validation} \\
\hline
Authentification sécurisée & Fonctionnel + intégration & Blocage avant vérification email, contrôle des rôles et réponses \texttt{403}. \\
\hline
Soumission de candidature fiable & Fonctionnel & Rejet des fichiers invalides, gestion du doublon (\texttt{409}), persistance correcte en base. \\
\hline
Suivi RH du pipeline & Intégration & Changement de statut répercuté en base et propagation des notifications. \\
\hline
Planification d'entretien & Fonctionnel + intégration & Création/replanification avec contrôle des états métier et notifications associées. \\
\hline
Isolation multi-site & Sécurité + intégration & Interdiction d'accès aux données hors périmètre site. \\
\hline
\end{tabularx}
\end{table}

\subsection{Conformité aux besoins}

Le système est globalement conforme aux besoins définis dans l'analyse :
\begin{itemize}
    \item besoin de digitalisation du recrutement ;
    \item besoin de transparence pour les candidats ;
    \item besoin de structuration et de traçabilité côté RH ;
    \item besoin d'isolation des données entre les sites.
\end{itemize}

Cette conformité confirme que la solution répond au problème initial posé
par l'entreprise, dans le périmètre défini pour ce projet de fin d'études.

% ── Conclusion du chapitre ──────────────────────────────────────
\section*{Conclusion du chapitre 5}
\addcontentsline{toc}{section}{Conclusion du chapitre 5}

Ce chapitre a présenté la stratégie de test, les outils utilisés, la
conception des cas, l'exécution et la validation du système.
Les résultats obtenus confirment une solution fonctionnelle, cohérente et
conforme aux besoins principaux dans le périmètre du projet.

Cette étape de validation clôt la partie technique et prépare la conclusion
générale du rapport.
